\title{Eagle and Finch: RWKV with Matrix-Valued States and Dynamic Recurrence}

\begin{abstract}
We present Eagle (RWKV-5) and Finch (RWKV-6), sequence models improving upon the RWKV (RWKV-4) \citep{peng2023rwkv} architecture. Our architectural design advancements include multi-headed matrix-valued states and a dynamic recurrence mechanism that improve expressivity while maintaining the inference efficiency characteristics of RNNs. We introduce a new multilingual corpus with 1.12 trillion tokens and a fast tokenizer based on greedy matching for enhanced multilinguality. We trained four Eagle models, ranging from 0.46 to 7.5 billion parameters, and two Finch models with 1.6 and 3.1 billion parameters and find that they achieve competitive performance across a wide variety of benchmarks.

We release all our models on HuggingFace under the Apache 2.0 license.\footnote{Models at: \url{https://huggingface.co/RWKV}\\Training code at: \url{https://github.com/RWKV/RWKV-LM}\\
Inference code at: \url{https://github.com/RWKV/ChatRWKV}\\
Time-parallel training code at: \url{https://github.com/RWKV/RWKV-infctx-trainer}}
\end{abstract}

\section{Introduction} \label{sec:introduction}

Advancements in Large Language Models (LLMs) have significantly impacted Natural Language Processing (NLP) tasks. The field has traditionally been dominated by the transfomer architecture \citep{vaswani2023attention}. However, the expressive attention mechanism of transformers leads them to suffer from quadratic time complexity with respect to input sequence length. Various methods have been proposed to achieve sub-quadratic time complexity without significantly changing the core attention mechanism, typically relying on some form of sparsity techniques~\citep{child2019generatingsparse, beltagy2020longformer, zaheer2020bigbird}.

Recent works have achieved sub-quadratic time complexity without significantly sacrificing performance by introducing new mechanisms to replace attention at the core of the Transformer architecture. These models include gated recurrences ~\citep{fu2023hungry,gu2023mamba,gu2021s4,sun2023retentive,katsch2023gateloop, qin2023hierarchically, smith2023s5}, gated convolutions~\citep{poli2023hyena, peng2023rwkv}, data-dependent linear attention~\citep{yang2023gated, linearTrans2020inputDepend}, sparse attentions~\citep{tay2020sparse,child2019sparsetransformer,zaheer2020bigbird,blockwindow2019sparsetransformer} and their combinations \citep{de2024griffin, qin2024transnormerllm, qin2022deviltransnormer}. We build off RWKV-4 introduced in \cite{peng2023rwkv}, which provides efficient inference and training along with a parallelizable implementation compared to competing architectures as shown in Table \ref{tab:costs}.

In this paper, we introduce two new architectures: \textbf{Eagle} (RWKV-5) and \textbf{Finch} (RWKV-6). First, Eagle improves upon the architecture and learned decay schedule from RWKV-4~\citep{peng2023rwkv} through the use of expressive multi-headed matrix-valued states (as opposed to vector-valued states), a reformulated receptance, and an additional gating mechanism. 

Finch further improves the expressivity and flexibility of the architecture by introducing new data-dependent functions for both the time-mixing and token-shift modules, consisting of parameterized linear interpolations. Additionally, Finch proposes a novel use of the Low Rank Adaptation~\citep{hu2022lora} function to allow for trainable weight matrices to efficiently augment the learned data decay vectors in a context-dependent manner.

Finally, we introduce a new tokenizer, the RWKV World Tokenizer, and a new dataset, RWKV World v2 (1.12 trillion tokens), specially designed to improve performance on multilingual and code data.

Through extensive experimentation, we show that the Eagle and Finch models perform competitively, or improve upon existing models under a wide variety of sequence modeling domains and tasks. Specifically, we evaluate our trained models on commonly used English-only and multilingual text benchmarks, associative recall, music modeling, and vision-language benchmarks. Our experiments demonstrate that the advancements in Eagle and Finch provide significant progress towards developing more efficient AI models In summary, our main contributions are:

\begin{itemize}
    \item The Eagle (RWKV-5) and Finch (RWKV-6) RWKV architectures, which significantly improve over RWKV-4 on benchmarks for LLMs.
    \item The RWKV World Tokenizer which contains underrepresented languages' vocabulary and which performs fast tokenization with Trie-based greedy matching.

    \item The RWKV World v2 public dataset, comprised of 1.12 trillion tokens of publicly available multilingual data.
    \item Public release of four pre-trained Eagle models, scaling from 0.46 to 7.5 billion parameters, and two Finch models, with 1.6 and 3.1 billion parameters. Demonstrating that these novel architectures are competitive to transformers when trained using enough FLOPs to make meaningful scaling conclusions.
    \item A completely open training pipeline to enable interpretability and reproducibility of alternative-architecture LLMs (See Table \ref{tab:compare}).

\end{itemize}

\section{Background} \label{sec:background}

Eagle and Finch are RNNs based on a multi-headed hybridization of the RWKV-4 architecture and linear attention. We discuss related work and the evolution of these two architectures below, with a more detailed review given in Appendix ~\ref{app_related_work}.

Recurrent Neural Networks (RNNs) are well suited to provide inexpensive inference on sequence modelling tasks, typically operating in O(1) time complexity per step with respect to sequence length. They model sequences with time dependencies by generating a hidden state $h_t$ at each time step, which is fed back in at the next time step as a secondary input. Classic RNNs (e.g. LSTM~\citep{hochreiter1997long} and GRU~\citep{cho2014learning}) became widely used for sequence modelling, but are difficult to parallelize across the time dimension for training.

The Transformer architecture has enjoyed remarkable success in generative sequence modelling, and language modelling in particular~\citep{vaswani2023attention, radford2018improving}, providing SOTA performance across many tasks. However, the use of multi-headed dot-product self-attention (MHA) leads to a quadratic time complexity with respect to sequence length.  The deficiencies of classic RNNs and Transformers led to many attempts to develop architectures incorporating the best features of both in a single model, namely O(1) per token time complexity and fast highly parallelizable training.

Linear Attention~\citep{schmidhuber1992unnormedlinattn, katharopoulos2020lineartransformrers} replaces the numerator of MHA's $\mathrm{softmax}(QK^T)V$ with $\phi(Q)\phi(K)^\mathrm{T}V$, allowing a reordering of operations via associativity to $\phi(Q)(\phi(K)^\mathrm{T}V)$, where $\phi$ represents a non-negative feature-map function.  It can be computed as an RNN in $O(1)$ time per step by adding $\phi(K_i^T)V_i$ to a recurrent state at each time step $i$, or trained in parallel much like MHA. This accomplishes the main goals outlined above, but naive linear attention suffers from significantly reduced performance compared to MHA-based transformers. 

A modified form of linear attention, the Attention Free Transformer (AFT) \citep{zhai2021attention}, paved the way for the RWKV architecture, by using a number of attention heads equal to the size of the feature dimension and incorporating a set of learned pairwise positional biases, denoted as $w$.

\begin{align}
\mathrm{AFTAttn}_t &= \sigma_q(q_t) \odot \frac{\sum_{i=1}^{t} \exp(k_i + w_{i,t}) \odot v_i}{\sum_{i=1}^{t} \exp(k_i + w_{i,t})}
\end{align}

RWKV-4 reformulates the AFT equation by replacing the pair-wise positional biases with a channel-wise vector of additive weight decay rates $w$. It also adds a bonus term $u$ to offset the weight of only the current input specially.

\begin{align}
    \mathrm{wkv}_t &= \frac{ \sum_{i=1}^{t-1} \exp(-(t-1-i)w+k_i) \odot v_i + \exp(u+k_t) \odot v_t}{\sum_{i=1}^{t-1} \exp(-(t-1-i)w+k_i) + \exp(u+k_t)}.
\end{align}

RWKV-4 also adds token-shift and gating to both attention and feed-forward sub-blocks of transformer, and small embedding initialization and normalization to quickly arrive at well-distributed token embeddings. Combining all of these architectural changes led RWKV-4 to become the first RNN to rival the performance of Transformers, while maintaining fast parallelizable training and $O(1)$ time complexity per token.

There has been a recent revival of RNNs in NLP research \citep{tiezzi2024resurgence}. HGRN\citep{qin2023hierarchically} is a recent time-parallelizable data-dependent RNN that employs input and forget gates. TransNormer\citep{qin2022deviltransnormer} applies RMSNorm to linear attention to bound its output. Other new time-parallelizable data-dependent RNNs have also been invented concurrently with our work including GLA \citep{yang2023gated} and Griffin \citep{de2024griffin}. 

State Space Models (SSMs) employ a hidden state of basis function weights to model an approximation of the input function \citep{gu2020hippo}, updating that hidden state via a differential equation. Earlier SSMs \citep{gu2022efficiently} were historically computed using long convolutions in $O(N\log N)$ time per sequence, but could also be formulated as a recurrent network. Recently, it has been shown that SSMs can be parallelized across the time dimension via techniques including associative scan \citep{smith2023s5}.  A new class of SSMs has also emerged concurrently with our work \citep{katsch2023gateloop,gu2023mamba} that feature data-dependent $A$ and $B$ terms, which function similarly to the data-dependent dynamic recurrence used in Finch. 

\section{Eagle/Finch Architecture}\label{sec:rwkv_model}
We refine the RWKV architecture in two steps, and observe significant modeling improvements with each. Compared to the baseline RWKV-4, Eagle adds matrix-valued attention states, LayerNorm over the attention heads, SiLU attention gating, and improved initialization. It also removes the Sigmoid activation of receptance. Finch further applies data-dependence to the decay schedule and token-shift.

The core architecture remains similar to that of RWKV-4, consisting of a series of stacked residual blocks shaped like a traditional Transformer. Following notation from \citep{NEURIPS2021mixer}, each block contains one Pre-LayerNorm Time-Mixing sub-layer followed by one Pre-LayerNorm Channel-Mixing sub-layer, as depicted in \hyperref[fig:blocks]{Figure \ref*{fig:blocks}}, left. These correspond to the traditional Attention and Feed Forward Network sub-layers of the Transformer. See Appendix \ref{app_architecture} for more details on our training implementation and the differences from RWKV-4, and Section \ref{sec:Speed and Memory Benchmarks} for speed and memory benchmarks.

\section{Method}\label{sec:method}
In this section, we use $D$ to denote the model dimension, and unless explicitly stated, all vectors appearing in this section are dimension $D/h$, where $h$ denotes the number of heads, belonging to $\mathbb{R}^{(D/h)}$. For compactness and simplicity we show calculations per-head, eliding the head index. We use the convention that all vectors are row vectors unless explicitly transposed, so all matrices operate on the right side. We use the square subscript to denote a variable.

\subsection{Eagle}

\subsubsection{Eagle Token Shift}

We adopt the Token Shift technique from the previous RWKV, similar to a 1D causal convolution of size = 2, as can be seen in \hyperref[fig:blocks]{Figure \ref*{fig:blocks}}, center-bottom. To better introduce the Token Shift technique, we define some notation. The \textbf{l}inear int\textbf{erp}olation (lerp) between $x_t$ and $x_{t-1}$ used in RWKV-4 and Eagle Token Shift is defined as:

\begin{align}
\label{eq:lerp5}
\mathrm{lerp}_\square (a, b) &= a + (b-a) \odot \mu_\square
\end{align}

where each $\mu_\square\in\mathbb{R}^{D}$ is a learnable vector.

Token Shift allows the model to learn how much new versus old information should be allocated per time step to each channel of receptance, key, value, and gate vectors ($r$, $k$, $v$, and $g$ respectively) independently and uniquely for each head. This makes it possible to form induction heads \citep{elhage2021mathematical} within a single layer since even a single head can directly accumulate both past and current token data into separate subspaces within these vectors.

\subsubsection{Eagle Time Mixing}

The formula of Eagle Time Mixing can be written as follows:

\begin{align}\label{eq:token-shift5}
\square_{t} &= \mathrm{lerp}_{\square}(x_t, x_{t-1})\bm{W}_{\square}, \quad \square\in \{r,k,v,g\}\\
w &= \exp(-\exp(\omega))
\end{align}

\begin{align}
\label{eq:time-mix5}
\bm{wkv}_{t} &= \mathrm{diag}(u)\cdot k_{t}^\mathrm{T} \cdot v_{t} + \sum_{i=1}^{t-1} \mathrm{diag}(w)^{t-1-i} \cdot  k_{i}^\mathrm{T} \cdot v_{i}  \in \mathbb{R}^{(D/h) \times (D/h)}  \\
  \label{eq:time-mixer5out}
o_t &= \mathrm{concat}\left(\mathrm{SiLU}(g_t) \odot \layernorm(r_{t}  \cdot \bm{wkv}_{t})\right) \bm{W}_o \in \mathbb{R}^{D}
\end{align}

Where LayerNorm operates on each of $h$ heads separately, which is also equivalent to the GroupNorm (\cite{wu2018group}) operation on $h$ groups. It is also worth noting that $w$ is obtained from $w = \exp(-\exp(\omega))$, where $\omega \in \mathbb{R}^{D/h}$ are the actual headwise trainable parameters. This ensures that $w$ falls within the interval $(0,1)$, guaranteeing that $\mathrm{diag}(w)$ is a contraction matrix. 

The $\bm{wkv}_{t}$ attention calculation can alternatively be written in a recurrent form:

\begin{align}
\label{eq:time-mixer5}
\bm{wkv}' &= \bm{s} + \mathrm{diag}(u) \cdot k^\mathrm{T} \cdot v \\
\label{eq:time-mixer5state}
\bm{s}' &= \mathrm{diag}(w) \cdot \bm{s} + k^\mathrm{T} \cdot v
\end{align}

RWKV's $\bm{wkv}$ term can be considered a decay-based equivalent to the normalised $k^\mathrm{T}v$ term in Linear Attention. It is instructive to note how for a given head $j$ the recurrent state $s$ is a sum of $k^Tv$ where each channel of $s$ individually decays by the corresponding channel of $w$ at each time step. Prior to the application of the receptance vector, gating, and output weights, a per-channel learned boost $u$ is multiplied with the current token's $k^\mathrm{T}v$ and summed with the state, as can be seen in \hyperref[fig:blocks]{Figure \ref*{fig:blocks}}, top-right. This gives the current token special treatment relative to the sum of past tokens contained within the decaying state history. The receptance is multiplied by this sum, acting like the query term in Linear Attention.

\subsubsection{Channel Mixing} In both Eagle and Finch, the Channel Mixing module is identical to the previous RWKV-4 architecture, except for a slightly reduced hidden dimension from $4D$ to $3.5D$. This reduction accounts for new gating weights in Eagle Time Mixing to ensure an equi-parameter relation with the prior model at the same number of layers and embedding dimension. We do not further reduce the hidden dimension in Finch despite adding a small number of new parameters for LoRA weights. The formulas for Channel Mixing are the same as RWKV-4, but we restate them here to ensure notational consistency, using linear interpolation from Equation~\ref{eq:lerp5}:

\begin{align}
\label{eq:channel-mix5}
r'_t &= \mathrm{lerp}_{r'} (x'_t, x'_{t-1}) \bm{W}_{r'}\in \mathbb{R}^{D}\\
k'_t &= \mathrm{lerp}_{k'} (x'_t, x'_{t-1}) \bm{W}_{k'}\in \mathbb{R}^{3.5D}\\
v'_t &= \mathrm{ReLU}(k'_t)^2 \bm{W}_{v'}\in \mathbb{R}^{D}\\ 
o'_t &= \sigma(r'_t) \odot v'_t\in \mathbb{R}^{D}
\end{align}

\subsection{Finch}

\subsubsection{Finch Token Shift}

The \textbf{d}ata-\textbf{d}ependent \textbf{l}inear int\textbf{erp}olation (ddlerp) between $x_t$ and $x_{t-1}$ used in Finch Token Shift is defined as:

\begin{align}
\mathrm{lora}_\square(x) &= \lambda_\square + \tanh(x\bm{A}_\square)\bm{B}_\square \\
\mathrm{ddlerp}_\square (a, b) &= a + (b-a) \odot \mathrm{lora}_\square(a + (b-a) \odot \mu_x)
\end{align}

where $\mu_x$ and each $\lambda_\square$ introduce a trainable vector of dimension $D$ and each $\bm{A}_\square \in \mathbb{R}^{D \times 32}$, $\bm{B}_\square \in \mathbb{R}^{32 \times D}$ introduce new trainable weight matrices. For the special case of $\mathrm{LoRA}_\omega$ seen below we introduce double-sized trainable weight matrices $\bm{A}_\omega \in \mathbb{R}^{D \times 64}$, $\bm{B}_\omega \in \mathbb{R}^{64 \times D}$. A schematic representation can be found in \hyperref[fig:blocks]{Figure \ref*{fig:blocks}}, bottom-right. Please note that future 7B and larger Finch models are expected to further increase the size of these weight matrices by double or more.

This new form of Token Shift enhanced with data-dependence is intended to expand the abilities of the model beyond the RWKV-4/Eagle style of Token Shift so that the amount of new and old data allocated per channel now depends on the input at both current and prior time steps. 

\subsubsection{Finch Time Mixing}

\begin{align}\label{eq:token-shift6}
\square_{t} &= \mathrm{ddlerp}_{\square}(x_t, x_{t-1})\bm{W}_{\square}, \quad \square\in \{r,k,v,g\}\\
d_t &= \mathrm{lora}_d( \mathrm{ddlerp}_d ( x_t, x_{t-1} ) )\\
w_t &= \exp(-\exp(d_t))
\end{align}

\begin{align}
\bm{wkv}_{t} &=  \mathrm{diag}(u)\cdot k_{t}^\mathrm{T} \cdot v_{t} + \sum_{i=1}^{t-1}  \mathrm{diag}\left(\bigodot_{j=i+1}^{t-1}w_{j}\right) \cdot  k_{i}^\mathrm{T} \cdot v_{i}  \in \mathbb{R}^{(D/h) \times (D/h)} \\
o_t &= \mathrm{concat}\left(\mathrm{SiLU}(g_t) \odot \layernorm(r_{t} \cdot \bm{wkv}_{t})\right)\bm{W}_o \in \mathbb{R}^{D} 
\end{align}

The $\bm{wkv}_{t}$ attention calculation can alternatively be written in a recurrent manner:

\begin{align}
\label{eq:time-mixer6}
\bm{wkv}' &= \bm{s} + \mathrm{diag}(u) \cdot k^\mathrm{T} \cdot v \\
\bm{s}' &= \mathrm{diag}(w) \cdot \bm{s} + k^\mathrm{T} \cdot v
\end{align}

Unlike in Eagle, $w_t$ here is not static across the sequence (dashed arrows in \hyperref[fig:blocks]{Figure \ref*{fig:blocks}}, left and top-right.). This is the core change to decay in Finch, as each channel of $w_t$ can now vary independently over time, in a data-dependent manner, whereas previously it was a fixed learned vector.

The new LoRA mechanisms above are used to take learned vectors, as seen in Eagle, and inexpensively augment them with additional offsets determined by the incoming input. Note that the LoRA process itself uses an Eagle style Token-Shifted value as its input, not just the latest token. The new time-varying decay $w_t$ goes one step further, applying LoRA again afterward. Intuitively, this is a second-order variant of Token-Shifting, allowing each channel of $w_t$ to vary based on a mix of the current and prior tokens, with the mix itself determined by aspects of both tokens.

\section{RWKV World Tokenizer}\label{sec:tokenizer}
Tokenization is important in language modelling as it conditions the learning relationships between tokens and the generation of new text based on those patterns. The numbers of tokens to build a single semantic chunk are, however, often very unequally distributed against non-European and other underrepresented languages. Byte-pair-encoding (BPE) based tokenizers which are trained with this inequality result in not only lower performances against underrepresented languages but also undue economic costs such as inference~\cite{ahia2023languageCost} and continual pre-training  with extended vocabulary~\cite{mala500lin2024, elyzallama2023}. To address these problems, we manually select tokens from multiple vocabulary files such that non-European languages are well represented.

To construct the tokenizer's vocabulary, we merge the vocabularies of the following tokenizers and then manually select the tokens for non-European languages.

\begin{itemize}
    \item \textbf{GPT-NeoX-20B \citep{black-etal-2022-gpt}:} \url{https://huggingface.co/EleutherAI/gpt-neox-20b}
    \item \textbf{GPT2 \citep{radford2019language}:} \url{https://huggingface.co/openai-community/gpt2}
    \item \textbf{cl100k\_base of tiktoken:} \url{https://github.com/openai/tiktoken}
    \item \textbf{Llama2 \citep{touvron2023llama}:} \url{https://huggingface.co/meta-llama/Llama-2-7b-hf}
    \item \textbf{Bloom \citep{workshop2023bloom}:} \url{https://huggingface.co/bigscience/bloom}
\end{itemize}

This tokenizer has a vocabulary size of $V=65536$, numbered from 0 through 65535, where tokens are arranged by their lengths in bytes. Below is a brief overview:

\begin{itemize}
    \item \textbf{Token 0:} Represents the boundary between text documents, known as \texttt{<EOS>} or \texttt{<SOS>}. This token doesn't encode any specific content and is only used for document separation.
    \item \textbf{Tokens 1-256:} Consist of byte encodings (Token $k$ encodes byte $k-1$), wherein tokens 1-128 correspond to standard ASCII characters.
    \item \textbf{Tokens 257-65529:} Tokens with a minimum length of 2 bytes in UTF-8, including words, prefixes and suffixes, accented letters, Chinese characters, Hangul, Hiragana, Katakana and emojis. For example, Chinese characters are allocated from token 10250 to 18493.
    \item \textbf{Token 65530-65535:} Reserved tokens for future use.
\end{itemize}

These designations are intended to enhance the tokenizer's efficiency on the multilingual corpus, as well as on source code of programming languages.

This tokenizer is implemented via a Trie (Prefix Tree) to boost speed while maintaining simplicity. Encoding is performed as matching the longest element in vocabulary with an input string from left to right. We note that our tokenizer's vocabulary construction is to mitigate \textit{undue} burden, which naive BPE and related methods cause, on minor languages.

\section{RWKV World v2 Dataset}\label{sec:dataset}

We train our models on the new \textbf{RWKV World v2 Dataset}, a new multilingual 1.12 trillion token dataset drawn from a wide variety of hand selected publicly available data sources. This dataset is designed to go beyond the English-heavy focus of many datasets widely used to train LLMs today. We do this to support usage by the majority of the worldwide population who are not native English speakers, to improve representation within model responses, and also to enable transfer learning so that our models can apply knowledge across cultures and locales. We put a strong emphasis on factual knowledge and code, but also on cultural works including stories, books, subtitles, and conversations. The source data is approximately 70\% English, 15\% multilingual, and 15\% code. We describe the components of our dataset in detail in
\hyperref[sec:training_dataset_details]{Appendix \ref*{sec:training_dataset_details}}.

\section{Pre-Trained Models}
\label{sec:pretrained_models}

We have pre-trained and publicly released the six Apache 2.0 licensed Eagle and Finch models: \textbf{Eagle 0.4B}, \textbf{Eagle 1.5B}, \textbf{Eagle 3B}, \textbf{Eagle 7B}, \textbf{Finch 1.6B}, and \textbf{Finch 3B}. All of the models were trained on the 1.12 trillion token RWKV World v2 multilingual corpus. See \hyperref[flop_count]{Appendix \ref*{flop_count}} for detailed parameter counts and FLOPs calculations.

\section{Language Modeling Experiments}\label{evaluations}

\subsection{LM Evaluation Harness Benchmarks}\label{subsec:evals}

To assess the performance of Eagle and Finch models, we evaluate on a series of common multi-lingual and English-focused benchmarks using lm\_evaluation\_harness \citep{gao10256836framework} as shown in Tables~\ref{tab:multilang_bench} and~\ref{tab:eng_bench}. We find that Eagle and Finch demonstrate exceptionally high capabilities on multi-lingual benchmarks, with nearly all results significantly outperforming the other similarly sized models we tested.

In figures \ref{fig:multilang-flops} and \ref{fig:eng-flops} we plot the accuracy versus FLOPs used to train various open models across a similar set of common benchmarks. For multilingual benchmarks, Eagle and Finch represent a substantial improvement to the Pareto frontier, achieving far higher scores than other models trained for a similar number of FLOPs. The two models additionally obtain competitive performance across these English benchmarks.

\subsection{Associative Recall}
\label{subsec:associative_recall}

Associative recall (AR) is a synthetic task designed to mimic the way that humans associate and retrieve information. It measures a model's proficiency in recalling information that was previously mentioned in context. Prior research suggests that a model's ability to perform AR is indicative of its effectiveness in in-context learning \citep{elhage2021mathematical, olsson2022incontext}. As a result, AR has been adopted as a benchmark in developing new language model architectural designs. \citep{fu2023hungry, poli2023hyena, lutati2023focus}. \citet{arora2023zoology} benchmarked a range of models for multi-query associative recall (MQAR) and identified a performance gap between various linear transformer architectures and the transformer with attention. In MQAR tasks, prior RWKV models demonstrated a correlation between model dimension and sequence length. To compare architectures, we trained models using RWKV-4, Eagle and Finch on MQAR, using identical criteria with various model dimensions and sequence lengths. Our findings reveal significant improvements in MQAR with Eagle and Finch. Notably, Finch achieves extremely high accuracy in MQAR in our tests, and outperforms all well-known non-transformer architectures previously used to train large language models. Our experiments reveal performance disparities between Mamba~\citep{gu2023mamba} and Finch, despite their shared architectural features such as matrix-valued state and data-dependent memory modification, suggesting different combinations of these elements result in superior performance.

\subsection{Long Context Experiments}\label{subsec:long_context_experiments}

We test loss versus sequence position on the PG19 \citep{raecompressive2019} test set of books from token 2048 onward across RWKV-4, Eagle, and Finch. We find that Eagle improves dramatically over RWKV-4 on this long sequence task, despite having been trained solely on sequence length 4096. Finch further improves on this test beyond Eagle, with loss continuing to drop further into the sequence. See \hyperref[fig:rwkv-ctxlen]{Figure \ref*{fig:rwkv-ctxlen}} for details.

\subsection{Bamboo Benchmark} \label{subsec:bamboo}
The Bamboo benchmark~\citep{dong2023bamboo} evaluates the overall long-context language modeling capability of LLMs from five aspects: question answering, hallucination detection, text sorting, language modeling, and code completion, comprising a total of ten evaluation tasks. We test models on the 4k version of the benchmark, which includes all ten tasks with a maximum context window length of 4k. We choose not to present results on the code completion task since all tested models failed to generate correct code completions for this task. In Table~\ref{tab:bamboo_bench}, we present the results of nine tasks, with either accuracy or F1 score, along with their average scores. At both the 1.5b and 3b scales, the latest Finch and Eagle models outperform the vanilla Mamba by at least a 7\% average score, while remaining comparable with the Mamba trained on Hermes data (\textit{i.e.}, only a 0.7\% drop in the average score). Note that, despite being trained on only 1.1T tokens, Eagle-7b consistently outperforms Pythia by an average of 13.5\% at the 7b scale, and it also surpasses LLaMA2-Chat-7b on several tasks in the Bamboo benchmark. These results demonstrate the superior capacity of the proposed Finch and Eagle models on a vast range of long-context tasks.

\section{Speed and Memory Benchmarks}
\label{sec:Speed and Memory Benchmarks}

We compare the speed and memory utilization of the Attention-like kernels for Finch, Mamba\footnote{We also plot Mamba 2x which uses 2 runs through the Mamba kernel instead of one. This is done to mimic the usage of twice the number of layers in Mamba vs Finch and Transformers}, and Flash Attention\footnote{We use the PyTorch Implementation of Flash Attention v2}~\citep{dao2023flashattention2} in Figures \ref{fig:memory_vs_sequence} and \ref{fig:time_vs_sequence}. For all benchmarks, we use a batch size of 8, a model dimension of 4096, and a head size of 64 for both Flash Attention and Finch. For Mamba, we employ a state dimension of 16, a model dimension of 8192, to mimic Mamba's usage of an expansion factor of 2. Our findings indicate that Finch's speed in training scales linearly with respect to sequence length, exhibiting similar scaling to Mamba. We find Finch is significantly faster than Flash Attention for sequence lengths beyond 4k, being around 4.2x faster for a sequence length of 16k. Furthermore, Finch consistently outperforms Mamba and Flash Attention in terms of memory usage, using 40\% and 17\% less memory usage than Flash Attention and Mamba respectively. Further optimization of our Finch CUDA implementation, including algorithmic improvements, are possible, and could lead to speed increases and greater parallelization. However, this optimization is left for future work.

\section{Multimodal Experiments}
\label{sec:multimodal_experiments}

In this section, we explore the capabilities of Eagle when extended to handle multimodal tasks, where the model processes and integrates textual inputs with inputs in a different domain.

\subsection{RWKV Music Modelling}
\label{sec:music-modelling}
To investigate the Eagle architecture's applicability to music modeling, we use the Irishman ABC music sheet dataset \citep{wu2023tunesformer} to train a new RWKV-5-Music model using the same hyperparameters as the existing RWKV-4-Music model. The loss of RWKV-5 is approximately 2\% lower than that of the previous generation model, and this improvement is primarily observed in the musical score part, indicating that RWKV-5 possesses stronger modeling and generalization capabilities than its predecessor. The model has a total of $L=24$ layers, with a dimension of $D=512$ and uses a byte-level tokenizer with $V=128$ tokens. The training context length is $1024$ bytes. We use all 2,162 pieces of music in the validation set and calculate the loss for each position from the start. The loss is averaged across all pieces of music, then Gaussian smoothed over the position in the sequence.

The figure \ref{fig:music} shows the loss as a function of position. Note that the first 30-100 bytes of the ABC format are the file header and control codes, followed by the musical scores. The loss of RWKV-5 is approximately 2\% lower than the previous generation model, and it is shown mainly in the musical score part, indicating that RWKV-5 has stronger modelling and generalization capabilities than its precedent model.

\subsection{VisualRWKV}
\label{sec:visualrwkv}

VisualRWKV is the visual-enhanced version of the RWKV language model, enabling RWKV to handle various visual tasks. Our VisualRWKV follows a similar architecture to popular vision-language models \citep{liu2023improvedllava15}. We present the architecture in Figure \ref{fig:visualrwkv}. It consists of a vision encoder and a language model. Specifically, we use CLIP \citep{radford2021learningclip} as the vision encoder and Eagle 1.5B and 3B as the language model. We use LLaVA-1.5 dataset ~\citep{liu2023improvedllava15}. To adapt Eagle to this multimodal task, we employ a two-stage instruction-tuning process to enhance model performance. Initially, we conduct pre-training for feature alignment, during which only the projection layer is subjected to updates, while the rest of the model is kept in a frozen state. Following this, we move on to the fine-tuning end-to-end stage, where both the projection layer and the RWKV language model are fine-tuned, and the vision encoder continue to be kept frozen. As shown in Table \ref{tab:visualrwkv_results}, we demonstrate that VisualRWKV's architecture is powerful for visual understanding and reasoning. With a smaller vision encoder CLIP-L (0.4B) and modest-sized LLMs of 1.5B and 3B, it achieves results comparable to the combination of CLIP-G (1.0B) and CLIP-H (1.0B) with larger LLMs of 7B and 13B. Moreover, in some benchmarks, it even outperforms larger models.

\section{RWKV on Audio}
\label{sec:audiorwkv}

AudioRWKV is the audio-specific version of RWKV, with a better process of the input audio spectrogram. Inspired by the VRWKV \citep{wang2024state}, we introduce a quad-directional shift (Q-Shift) to capture the neighboring relationships in two-dimensional audio spectrograms in the first step of each spatial-mix and channel-mix module. Specifically, the Q-Shift operation allows all tokens to be shifted and linearly interpolated with their neighboring tokens. We conduct experiments on the AudioSet \citep{gemmeke2017audio} dataset with various model sizes from 8.7M to 105M. As shown in Table \ref{tab:audiorwkv_results}, AudioRWKV-Tiny achieves a comparable performance with AST-AT by a smaller model size.

\section{Conclusions} \label{conclusions}
In this work, we introduced Eagle (RWKV-5) and Finch (RWKV-6), marking substantial progress in RNN-based language models by integrating multiheaded matrix-valued states and dynamic data-driven recurrence mechanisms. These models demonstrate exceptional performance on MQAR and diverse linguistic benchmarks, challenging the dominance of traditional Transformer architectures while retaining key RNN advantages. With models publicly available under the Apache 2.0 license and trained on an extensive multilingual corpus, our work not only advances the capabilities of language models but also emphasizes community accessibility and applicability across various domains. While acknowledging the computational and ethical challenges ahead, we hope that Eagle and Finch's efficient new architecture and wide availability will help push the boundaries of language modeling and pave the way for future innovations.

\paragraph{Limitations}
The Eagle and Finch models fall short on certain aspects that can be mitigated and addressed in future work.

We experimented with using Eagle as an embedding model on the Massive Text Embedding Benchmark (MTEB) \citep{muennighoff2023mteb} but were not able to get strong embedding performance. We believe that its state is a very high-quality embedding of the context but an appropriate method is required to aggregate the information content. We leave this to future work.

Because our training corpus contains some synthetic data from GPT-3.5 and ChatGPT, our released models exhibit behaviors similar to ChatGPT and will mimic ChatGPT's conversation style and tone. For instance, the model might occasionally claim that it is trained by OpenAI. However, this is not a general property the RWKV architecture but rather a specific outcome of the data and training process.

\paragraph{Future Work} \label{future_work}

Our 1.12 trillion token multilingual training corpus is much smaller than the training data sizes for contemporary models such as LLaMA2 \citep{touvron2023llama}, and expanding our training corpus to be more diverse and expansive is a key priority to improving model performance~\citep{albalak2024survey}. We also plan to train and release larger versions of Finch such as 7B and 14B parameters, and further extend its performance with reduced inference and training costs via Mixture of Experts~\citep{shazeer2017outrageously}.

\subsubsection*{Acknowledgments}
We thank Stability AI for the compute used to train our models and for technical support in the development of RWKV. We also thank the members of the RWKV and EleutherAI Discord servers for their help and work on further extending the applicability of RWKV to different domains. We also thank Shenzhen Yuanshi Intelligence Co., Ltd. for its contribution to the promotion and commercialization of RWKV. We thank Songlin Yang for assistance with the code and ideas for our time-parallel implementations.

\end{document}
